\documentclass[11pt,a4paper]{article}

% --------------------------------------------------
% Packages
% --------------------------------------------------
\usepackage[margin=1in]{geometry}
\usepackage{parskip}
\usepackage{booktabs}
\usepackage{array}
\usepackage{enumitem}
\usepackage{xcolor}
\usepackage{hyperref}
\usepackage{listings}
\usepackage{fancyhdr}
\usepackage{titlesec}
\usepackage{longtable}
\usepackage{amsmath}
\usepackage{graphicx}
\usepackage{microtype}

% --------------------------------------------------
% Colors
% --------------------------------------------------
\definecolor{darkblue}{RGB}{31,78,121}
\definecolor{lightblue}{RGB}{230,240,250}
\definecolor{codegray}{RGB}{245,245,245}
\definecolor{codegreen}{RGB}{40,120,70}
\definecolor{codered}{RGB}{170,50,50}

% --------------------------------------------------
% Hyperlinks
% --------------------------------------------------
\hypersetup{
    colorlinks=true,
    linkcolor=darkblue,
    urlcolor=darkblue,
    citecolor=darkblue
}

% --------------------------------------------------
% Section formatting
% --------------------------------------------------
\titleformat{\section}
  {\Large\bfseries\color{darkblue}}
  {\thesection.}
  {0.6em}
  {}

\titleformat{\subsection}
  {\large\bfseries}
  {\thesubsection.}
  {0.6em}
  {}

% --------------------------------------------------
% Header / footer
% --------------------------------------------------
\pagestyle{fancy}
\fancyhf{}
\lhead{GridCast}
\rhead{7-Day MLOps Project Plan}
\cfoot{\thepage}

% --------------------------------------------------
% Lists
% --------------------------------------------------
\setlist[itemize]{topsep=4pt,itemsep=2pt}
\setlist[enumerate]{topsep=4pt,itemsep=3pt}

% --------------------------------------------------
% Code formatting
% --------------------------------------------------
\lstset{
    backgroundcolor=\color{codegray},
    basicstyle=\ttfamily\small,
    breaklines=true,
    frame=single,
    rulecolor=\color{gray!35},
    columns=fullflexible,
    showstringspaces=false,
    keywordstyle=\color{darkblue}\bfseries,
    commentstyle=\color{codegreen},
    stringstyle=\color{codered},
    tabsize=4
}

% --------------------------------------------------
% Commands
% --------------------------------------------------
\newcommand{\done}{\textbf{\textcolor{codegreen}{✓}}}
\newcommand{\projectname}{GridCast}

% --------------------------------------------------
% Document
% --------------------------------------------------
\begin{document}

\begin{center}
    {\Huge \bfseries \projectname}\\[0.3cm]
    {\Large Production Energy Consumption Forecasting}\\[0.15cm]
    {\large Drift Detection, Automated Retraining, and Model Promotion}\\[0.6cm]

    \textbf{Project Duration:} 7 Days\\
    \textbf{Primary Goal:} MLOps Portfolio Project
\end{center}

\vspace{0.5cm}

\hrule

\vspace{0.5cm}

\section{Project Overview}

\projectname{} is an end-to-end MLOps portfolio project for forecasting electricity consumption.

The system forecasts the next 24 hours of aggregate electricity demand using historical consumption data. The project emphasizes the complete machine-learning lifecycle rather than advanced forecasting architectures.

The final system will demonstrate:

\begin{itemize}
    \item Reproducible data preprocessing
    \item Time-aware feature engineering
    \item Baseline forecasting
    \item XGBoost-based demand forecasting
    \item Experiment tracking with MLflow
    \item Model versioning and registry
    \item Data and pipeline versioning using DVC
    \item Containerized prediction serving
    \item Production-data simulation
    \item Forecast performance monitoring
    \item Data drift detection
    \item Automated model retraining
    \item Candidate-versus-production evaluation
    \item Model promotion and rollback
    \item Automated testing
    \item Continuous integration with GitHub Actions
\end{itemize}

The intended lifecycle is:

\begin{verbatim}
Data
  |
  v
Preprocessing
  |
  v
Feature Engineering
  |
  v
Model Training
  |
  v
MLflow Registry
  |
  v
Production Forecasting
  |
  +------------------+
  |                  |
  v                  v
Predictions        Actuals
  |                  |
  +--------+---------+
           |
           v
      Monitoring
       /      \
      v        v
   Drift    Performance
      \        /
       \      /
        v    v
     Retraining Gate
           |
           v
     Candidate Model
           |
           v
      Quality Gate
        /      \
       v        v
    Reject    Promote
\end{verbatim}

\section{Dataset}

Use the \textbf{UCI ElectricityLoadDiagrams20112014} dataset.

It contains electricity consumption measurements for hundreds of electricity clients and provides multiple years of historical consumption data.

For this project, the individual customer series will be aggregated into one system-level electricity demand time series.

The project will work with an hourly dataset of the form:

\begin{lstlisting}
timestamp               load
2012-01-01 00:00       14582.1
2012-01-01 01:00       13994.7
2012-01-01 02:00       13211.3
...
\end{lstlisting}

The target variable is:

\[
y_t = \text{aggregate electricity demand at hour } t
\]

The production objective is:

\[
\hat{y}_{t+1}, \hat{y}_{t+2}, \dots, \hat{y}_{t+24}
\]

where the system generates a 24-hour demand forecast.

\section{Technology Stack}

\begin{table}[h]
\centering
\renewcommand{\arraystretch}{1.2}
\begin{tabular}{>{\bfseries}p{0.31\textwidth} p{0.55\textwidth}}
\toprule
Component & Technology \\
\midrule
Programming Language & Python \\
Data Processing & Pandas, NumPy \\
Forecast Model & XGBoost \\
Baseline Models & Pandas / NumPy \\
Metrics & MAE, RMSE, MAPE \\
Experiment Tracking & MLflow \\
Model Registry & MLflow \\
Data Versioning & DVC \\
Pipeline Versioning & DVC \\
Drift Monitoring & Evidently \\
API & FastAPI \\
Validation & Pydantic \\
Testing & pytest \\
Containerization & Docker \\
CI & GitHub Actions \\
Code Quality & Ruff \\
Visualization & Matplotlib, MLflow, Evidently \\
\bottomrule
\end{tabular}
\end{table}

\section{Repository Structure}

The final repository should resemble:

\begin{lstlisting}
gridcast/
|
+-- data/
|   +-- raw/
|   +-- processed/
|   +-- production/
|
+-- notebooks/
|   +-- 01_data_exploration.ipynb
|
+-- src/
|   +-- data/
|   |   +-- preprocess.py
|   |   +-- validate.py
|   |
|   +-- features/
|   |   +-- build_features.py
|   |
|   +-- models/
|   |   +-- train.py
|   |   +-- evaluate.py
|   |   +-- predict.py
|   |
|   +-- monitoring/
|       +-- performance.py
|       +-- drift.py
|       +-- retrain.py
|
+-- api/
|   +-- main.py
|
+-- scripts/
|   +-- simulate_production.py
|   +-- run_demo.py
|
+-- tests/
|   +-- test_data.py
|   +-- test_features.py
|   +-- test_model.py
|   +-- test_monitoring.py
|   +-- test_registry.py
|   +-- test_api.py
|
+-- reports/
|   +-- drift/
|   +-- figures/
|
+-- models/
|
+-- .github/
|   +-- workflows/
|       +-- ci.yml
|
+-- Dockerfile
+-- dvc.yaml
+-- params.yaml
+-- requirements.txt
+-- Makefile
+-- README.md
\end{lstlisting}

\newpage

% ==================================================
% DAY 1
% ==================================================

\section{Day 1 --- Dataset, EDA, and Preprocessing}

\subsection{Objective}

The goal for Day 1 is to build a clean, validated, hourly electricity-demand dataset and establish the repository structure.

By the end of the day, the raw dataset should be transformable into the processed dataset using a reproducible script.

\subsection{1. Repository Setup}

Create the project directory structure.

Initialize Git:

\begin{lstlisting}[language=bash]
git init
\end{lstlisting}

Create a virtual environment:

\begin{lstlisting}[language=bash]
python -m venv .venv
source .venv/bin/activate
\end{lstlisting}

On Windows:

\begin{lstlisting}[language=bash]
.venv\Scripts\activate
\end{lstlisting}

Install the initial dependencies:

\begin{lstlisting}[language=bash]
pip install pandas numpy scikit-learn xgboost matplotlib \
mlflow dvc pytest ruff fastapi uvicorn evidently
\end{lstlisting}

Store dependencies:

\begin{lstlisting}[language=bash]
pip freeze > requirements.txt
\end{lstlisting}

\subsection{2. Dataset Storage}

Store the original dataset inside:

\begin{lstlisting}
data/raw/
\end{lstlisting}

The raw file must never be manually modified.

Processed files should be written to:

\begin{lstlisting}
data/processed/
\end{lstlisting}

\subsection{3. Exploratory Data Analysis}

Create:

\begin{lstlisting}
notebooks/01_data_exploration.ipynb
\end{lstlisting}

Investigate:

\begin{itemize}
    \item Start and end timestamps
    \item Sampling frequency
    \item Number of customer time series
    \item Missing values
    \item Duplicate timestamps
    \item Zero-consumption periods
    \item Potential abnormal values
    \item Consumption seasonality
\end{itemize}

Generate the following plots:

\begin{enumerate}
    \item Aggregate load over the full dataset
    \item One week of electricity consumption
    \item One month of electricity consumption
    \item Average electricity consumption by hour
    \item Average electricity consumption by weekday
    \item Average electricity consumption by month
\end{enumerate}

EDA should remain lightweight. The primary purpose is to understand the data sufficiently to build the forecasting pipeline.

\subsection{4. Aggregate Customer Demand}

For each timestamp, aggregate electricity usage:

\[
L_t = \sum_{i=1}^{N} x_{i,t}
\]

where:

\begin{itemize}
    \item $x_{i,t}$ is the load of customer $i$ at time $t$
    \item $N$ is the number of customers
    \item $L_t$ is the system-wide electricity demand
\end{itemize}

\subsection{5. Hourly Resampling}

Convert 15-minute consumption measurements into hourly consumption.

The output should resemble:

\begin{lstlisting}
timestamp,load
2012-01-01 00:00:00,14582.1
2012-01-01 01:00:00,13994.7
2012-01-01 02:00:00,13211.3
\end{lstlisting}

Implement this in:

\begin{lstlisting}
src/data/preprocess.py
\end{lstlisting}

The preprocessing pipeline should perform:

\begin{verbatim}
Raw Data
   |
   v
Parse Timestamp
   |
   v
Convert Numeric Columns
   |
   v
Aggregate Customers
   |
   v
Resample to Hourly
   |
   v
Validate
   |
   v
Save Processed Dataset
\end{verbatim}

\subsection{6. Data Validation}

Create:

\begin{lstlisting}
src/data/validate.py
\end{lstlisting}

Implement checks such as:

\begin{lstlisting}[language=Python]
assert df["timestamp"].is_monotonic_increasing
assert not df["timestamp"].duplicated().any()
assert df["load"].isna().sum() == 0
assert (df["load"] >= 0).all()
\end{lstlisting}

Also verify:

\begin{itemize}
    \item Hourly timestamp frequency
    \item Required columns exist
    \item No unexpected infinite values
    \item Target values are numeric
\end{itemize}

\subsection{7. Define Time Splits}

Do not use a random train-test split.

Use chronological partitions.

Recommended split:

\begin{table}[h]
\centering
\begin{tabular}{lll}
\toprule
Dataset & Start & End \\
\midrule
Training & 2012-01-01 & 2013-06-30 \\
Validation & 2013-07-01 & 2013-12-31 \\
Production Simulation & 2014-01-01 & 2014-12-31 \\
\bottomrule
\end{tabular}
\end{table}

The entire 2014 period should remain untouched during initial model development.

\subsection{8. Initialize DVC}

Initialize DVC:

\begin{lstlisting}[language=bash]
dvc init
\end{lstlisting}

Track the raw dataset:

\begin{lstlisting}[language=bash]
dvc add data/raw/<dataset-file>
\end{lstlisting}

Commit the metadata:

\begin{lstlisting}[language=bash]
git add .
git commit -m "Add electricity dataset and preprocessing pipeline"
\end{lstlisting}

\subsection{9. Day 1 Tests}

Implement:

\begin{itemize}
    \item \texttt{test\_timestamp\_parsing}
    \item \texttt{test\_no\_duplicate\_timestamps}
    \item \texttt{test\_hourly\_resampling}
    \item \texttt{test\_non\_negative\_load}
    \item \texttt{test\_no\_missing\_target}
\end{itemize}

\subsection{Day 1 Deliverables}

\begin{itemize}
    \item[\done] GitHub repository
    \item[\done] Project directory structure
    \item[\done] Reproducible Python environment
    \item[\done] Raw dataset
    \item[\done] Hourly processed dataset
    \item[\done] Basic EDA
    \item[\done] Preprocessing script
    \item[\done] Data validation
    \item[\done] DVC initialization
    \item[\done] Initial tests
\end{itemize}

\textbf{Stop condition:} One command should regenerate the processed dataset from the raw data.

\newpage

% ==================================================
% DAY 2
% ==================================================

\section{Day 2 --- Feature Engineering and Forecast Model}

\subsection{Objective}

Build a reproducible forecasting pipeline containing baseline models and an XGBoost model.

The model should be useful, but model optimization should remain secondary to the MLOps components.

\subsection{1. Calendar Features}

Create:

\begin{lstlisting}
src/features/build_features.py
\end{lstlisting}

Generate:

\begin{itemize}
    \item Hour
    \item Day of week
    \item Day of month
    \item Month
    \item Week of year
    \item Weekend indicator
\end{itemize}

\subsection{2. Lag Features}

Generate lagged electricity-demand values:

\begin{lstlisting}
lag_1
lag_2
lag_3
lag_6
lag_12
lag_24
lag_48
lag_168
\end{lstlisting}

Important examples:

\[
\text{lag}_{24}(t) = y_{t-24}
\]

\[
\text{lag}_{168}(t) = y_{t-168}
\]

The 168-hour lag represents electricity demand at approximately the same hour one week earlier.

\subsection{3. Rolling Features}

Generate:

\begin{lstlisting}
rolling_mean_6
rolling_mean_24
rolling_mean_168

rolling_std_24
rolling_std_168

rolling_min_24
rolling_max_24
\end{lstlisting}

Prevent future-data leakage by shifting before calculating rolling statistics.

Correct approach:

\begin{lstlisting}[language=Python]
df["rolling_mean_24"] = (
    df["load"]
    .shift(1)
    .rolling(24)
    .mean()
)
\end{lstlisting}

Avoid centered rolling windows.

\subsection{4. Baseline Models}

Implement two baseline forecasts.

\subsubsection{Yesterday Baseline}

\[
\hat{y}_t = y_{t-24}
\]

\subsubsection{Weekly Seasonal Baseline}

\[
\hat{y}_t = y_{t-168}
\]

Evaluate both using:

\begin{itemize}
    \item Mean Absolute Error
    \item Root Mean Squared Error
    \item Mean Absolute Percentage Error
\end{itemize}

\subsection{5. Evaluation Metrics}

Mean Absolute Error:

\[
MAE =
\frac{1}{n}
\sum_{i=1}^{n}
|y_i-\hat{y}_i|
\]

Root Mean Squared Error:

\[
RMSE =
\sqrt{
\frac{1}{n}
\sum_{i=1}^{n}
(y_i-\hat{y}_i)^2
}
\]

Mean Absolute Percentage Error:

\[
MAPE =
\frac{100}{n}
\sum_{i=1}^{n}
\left|
\frac{y_i-\hat{y}_i}{y_i}
\right|
\]

\subsection{6. Train XGBoost}

Create:

\begin{lstlisting}
src/models/train.py
\end{lstlisting}

A reasonable starting model is:

\begin{lstlisting}[language=Python]
from xgboost import XGBRegressor

model = XGBRegressor(
    n_estimators=500,
    max_depth=6,
    learning_rate=0.05,
    subsample=0.8,
    colsample_bytree=0.8,
    objective="reg:squarederror"
)
\end{lstlisting}

Avoid extensive hyperparameter optimization during the first week.

\subsection{7. Model Evaluation}

Create:

\begin{lstlisting}
src/models/evaluate.py
\end{lstlisting}

Save metrics to:

\begin{lstlisting}
reports/model_metrics.json
\end{lstlisting}

Example:

\begin{lstlisting}
{
    "mae": 421.7,
    "rmse": 601.4,
    "mape": 0.048
}
\end{lstlisting}

\subsection{8. Generate Forecast Visualizations}

Create:

\begin{itemize}
    \item Actual versus predicted load for one week
    \item Actual versus predicted load for one month
    \item Residual distribution
    \item Forecast error by hour
    \item Baseline versus XGBoost comparison
\end{itemize}

\subsection{9. Twenty-Four-Hour Forecast Output}

The system should generate a forecast table similar to:

\begin{lstlisting}
timestamp               prediction
2014-01-01 00:00        18431
2014-01-01 01:00        17902
2014-01-01 02:00        17320
...
2014-01-01 23:00        20741
\end{lstlisting}

All features for each prediction must use information that would have been available at forecast time.

\subsection{10. Day 2 Tests}

Implement tests for:

\begin{itemize}
    \item Feature-column existence
    \item Correct lag calculation
    \item Correct rolling statistics
    \item No future leakage
    \item Model-training completion
    \item Prediction shape
    \item Metric generation
    \item Basic performance sanity threshold
\end{itemize}

Example:

\begin{lstlisting}[language=Python]
assert model_mape < 0.20
\end{lstlisting}

\subsection{Day 2 Deliverables}

\begin{itemize}
    \item[\done] Feature pipeline
    \item[\done] Yesterday baseline
    \item[\done] Weekly baseline
    \item[\done] XGBoost model
    \item[\done] MAE, RMSE, and MAPE
    \item[\done] Forecast plots
    \item[\done] Serialized model
    \item[\done] Model tests
\end{itemize}

\textbf{Stop condition:} One command should train the model and produce measurable forecasts.

\newpage

% ==================================================
% DAY 3
% ==================================================

\section{Day 3 --- MLflow, Model Registry, and DVC Pipeline}

\subsection{Objective}

Turn the forecasting implementation into a reproducible machine-learning lifecycle.

\subsection{1. Configure MLflow}

Start a local MLflow server:

\begin{lstlisting}[language=bash]
mlflow server
\end{lstlisting}

Create an experiment called:

\begin{lstlisting}
gridcast-energy-forecasting
\end{lstlisting}

\subsection{2. Log Model Parameters}

Log:

\begin{itemize}
    \item Number of estimators
    \item Maximum depth
    \item Learning rate
    \item Subsample
    \item Column sampling
    \item Training start date
    \item Training end date
    \item Number of features
\end{itemize}

\subsection{3. Log Metrics}

Log:

\begin{itemize}
    \item Training MAE
    \item Validation MAE
    \item Validation RMSE
    \item Validation MAPE
\end{itemize}

\subsection{4. Log Artifacts}

Store:

\begin{itemize}
    \item Trained model
    \item Feature importance
    \item Forecast plots
    \item Metrics JSON
    \item Feature list
\end{itemize}

\subsection{5. Run Controlled Experiments}

Run approximately three to five meaningful experiments.

Example:

\begin{lstlisting}
Run 1: Base XGBoost
Run 2: Increased max_depth
Run 3: Additional rolling features
Run 4: Final candidate
\end{lstlisting}

Avoid dozens of arbitrary runs.

\subsection{6. Model Registry}

Register the best model as:

\begin{lstlisting}
GridCastForecaster
\end{lstlisting}

The logical model lifecycle should be:

\begin{verbatim}
Version 1
   |
   v
Candidate
   |
   v
Validation Gate
   |
   v
Production
\end{verbatim}

\subsection{7. Model Metadata}

Store metadata such as:

\begin{lstlisting}
{
    "model_version": "1",
    "training_start": "2012-01-01",
    "training_end": "2013-06-30",
    "features": [
        "hour",
        "day_of_week",
        "lag_24",
        "lag_168"
    ],
    "validation_mape": 0.049
}
\end{lstlisting}

\subsection{8. DVC Pipeline}

Create a \texttt{dvc.yaml} pipeline containing:

\begin{verbatim}
preprocess
   |
   v
features
   |
   v
train
   |
   v
evaluate
\end{verbatim}

Example:

\begin{lstlisting}
stages:

  preprocess:
    cmd: python -m src.data.preprocess
    deps:
      - data/raw/
      - src/data/preprocess.py
    outs:
      - data/processed/hourly_load.csv

  features:
    cmd: python -m src.features.build_features
    deps:
      - data/processed/hourly_load.csv
      - src/features/build_features.py
    outs:
      - data/processed/features.parquet

  train:
    cmd: python -m src.models.train
    deps:
      - data/processed/features.parquet
      - src/models/train.py
      - params.yaml
    outs:
      - models/model.pkl

  evaluate:
    cmd: python -m src.models.evaluate
    deps:
      - models/model.pkl
    metrics:
      - reports/model_metrics.json
\end{lstlisting}

Running:

\begin{lstlisting}[language=bash]
dvc repro
\end{lstlisting}

should reproduce the ML pipeline.

\subsection{9. Configuration}

Create:

\begin{lstlisting}
params.yaml
\end{lstlisting}

Example:

\begin{lstlisting}
model:
  n_estimators: 500
  max_depth: 6
  learning_rate: 0.05

features:
  lags:
    - 1
    - 24
    - 168

  rolling_windows:
    - 24
    - 168
\end{lstlisting}

\subsection{Day 3 Deliverables}

\begin{itemize}
    \item[\done] MLflow tracking
    \item[\done] Multiple experiment runs
    \item[\done] Logged artifacts
    \item[\done] Registered model
    \item[\done] DVC training pipeline
    \item[\done] Configuration file
    \item[\done] Model metadata
\end{itemize}

\textbf{Stop condition:} The project should be reproducible through \texttt{dvc repro}.

\newpage

% ==================================================
% DAY 4
% ==================================================

\section{Day 4 --- FastAPI, Model Serving, and Docker}

\subsection{Objective}

Expose the production model through a small service and containerize it.

\subsection{1. FastAPI Application}

Create:

\begin{lstlisting}
api/main.py
\end{lstlisting}

Implement:

\begin{lstlisting}
GET  /health
GET  /model-info
POST /forecast
\end{lstlisting}

\subsection{2. Health Endpoint}

Expected response:

\begin{lstlisting}
{
    "status": "healthy"
}
\end{lstlisting}

\subsection{3. Model Information Endpoint}

Example:

\begin{lstlisting}
{
    "model": "GridCastForecaster",
    "version": "2",
    "algorithm": "XGBoost",
    "training_end": "2013-12-31",
    "validation_mape": 0.048
}
\end{lstlisting}

This demonstrates production model traceability.

\subsection{4. Forecast Endpoint}

Recommended input:

\begin{lstlisting}
{
    "forecast_date": "2014-05-01"
}
\end{lstlisting}

Example output:

\begin{lstlisting}
{
    "model_version": "2",
    "forecast_date": "2014-05-01",
    "predictions": [
        {
            "hour": "00:00",
            "predicted_load": 18240.3
        },
        {
            "hour": "01:00",
            "predicted_load": 17681.9
        }
    ]
}
\end{lstlisting}

Feature construction should remain internal to the service.

\subsection{5. Input Validation}

Use Pydantic to reject:

\begin{itemize}
    \item Invalid timestamps
    \item Dates outside the available simulation range
    \item Missing required fields
    \item Malformed request payloads
\end{itemize}

\subsection{6. Logging}

Log at minimum:

\begin{itemize}
    \item Request timestamp
    \item Model version
    \item Forecast date
    \item Number of predictions
    \item Request latency
    \item Error status
\end{itemize}

\subsection{7. Docker}

Create a Dockerfile similar to:

\begin{lstlisting}
FROM python:3.11-slim

WORKDIR /app

COPY requirements.txt .

RUN pip install --no-cache-dir -r requirements.txt

COPY . .

CMD ["uvicorn", "api.main:app",
     "--host", "0.0.0.0",
     "--port", "8000"]
\end{lstlisting}

Build:

\begin{lstlisting}[language=bash]
docker build -t gridcast .
\end{lstlisting}

Run:

\begin{lstlisting}[language=bash]
docker run -p 8000:8000 gridcast
\end{lstlisting}

Swagger documentation should then be available through:

\begin{lstlisting}
http://localhost:8000/docs
\end{lstlisting}

\subsection{8. API Tests}

Test:

\begin{itemize}
    \item \texttt{/health} returns HTTP 200
    \item \texttt{/model-info} returns metadata
    \item \texttt{/forecast} returns 24 predictions
    \item Invalid input returns a 4xx response
\end{itemize}

\subsection{Day 4 Deliverables}

\begin{itemize}
    \item[\done] FastAPI application
    \item[\done] Health endpoint
    \item[\done] Model information endpoint
    \item[\done] Forecast endpoint
    \item[\done] Request validation
    \item[\done] Logging
    \item[\done] Docker image
    \item[\done] API tests
\end{itemize}

\textbf{Stop condition:} A user can start the Docker container and obtain forecasts.

\newpage

% ==================================================
% DAY 5
% ==================================================

\section{Day 5 --- Production Simulation and Performance Monitoring}

\subsection{Objective}

Treat the held-out 2014 dataset as live production data arriving over time.

\subsection{1. Production Simulator}

Create:

\begin{lstlisting}
scripts/simulate_production.py
\end{lstlisting}

Replay production data in weekly batches.

Conceptually:

\begin{verbatim}
Week 1
  |
  v
Generate Forecasts
  |
  v
Actual Data Arrives
  |
  v
Calculate Metrics
  |
  v
Store Results
  |
  v
Week 2
  |
  v
Repeat
\end{verbatim}

\subsection{2. Prediction Log}

Store production forecasts in:

\begin{lstlisting}
data/production/predictions.csv
\end{lstlisting}

Recommended schema:

\begin{lstlisting}
timestamp,prediction,actual,model_version
2014-01-01 00:00,18301,18741,1
2014-01-01 01:00,17920,18120,1
\end{lstlisting}

\subsection{3. Performance Monitoring}

Create:

\begin{lstlisting}
src/monitoring/performance.py
\end{lstlisting}

Calculate rolling or weekly:

\begin{itemize}
    \item MAE
    \item RMSE
    \item MAPE
\end{itemize}

Save:

\begin{lstlisting}
reports/production_metrics.csv
\end{lstlisting}

Example:

\begin{lstlisting}
week,model_version,mae,rmse,mape
1,1,421,580,0.041
2,1,440,604,0.044
3,1,471,630,0.049
\end{lstlisting}

\subsection{4. Define Performance Degradation}

Do not trigger retraining based on one bad prediction.

Suppose validation MAPE is:

\[
MAPE_{\text{baseline}} = 5\%
\]

Possible monitoring levels:

\begin{itemize}
    \item Warning: rolling MAPE $> 6\%$
    \item Retraining candidate: rolling MAPE $> 7\%$
\end{itemize}

A more flexible rule is:

\[
MAPE_{\text{current}}
>
1.30 \times
MAPE_{\text{baseline}}
\]

In code:

\begin{lstlisting}[language=Python]
performance_degraded = (
    rolling_mape >
    baseline_mape * 1.30
)
\end{lstlisting}

\subsection{5. Monitoring Visualizations}

Generate:

\begin{itemize}
    \item Production MAPE over time
    \item Production MAE over time
    \item Actual versus forecast
    \item Model version over time
\end{itemize}

A useful plot should visibly show when model quality deteriorates and when a new model version enters production.

\subsection{Day 5 Deliverables}

\begin{itemize}
    \item[\done] Production simulation
    \item[\done] Prediction history
    \item[\done] Rolling metrics
    \item[\done] Performance thresholds
    \item[\done] Monitoring plots
    \item[\done] Model-version tracking
\end{itemize}

\textbf{Stop condition:} The production simulator should replay data and show forecast performance changing over time.

\newpage

% ==================================================
% DAY 6
% ==================================================

\section{Day 6 --- Drift Detection, Retraining, and Model Promotion}

\subsection{Objective}

Implement the most important MLOps lifecycle:

\begin{verbatim}
Monitor
   |
   v
Detect Drift
   |
   v
Detect Performance Degradation
   |
   v
Retrain Candidate
   |
   v
Evaluate Candidate
   |
   v
Promote or Reject
\end{verbatim}

\subsection{1. Drift Windows}

Define a reference dataset and a current dataset.

Example reference period:

\begin{lstlisting}
2013-07-01 to 2013-12-31
\end{lstlisting}

Current production window:

\begin{lstlisting}
Most recent 30 production days
\end{lstlisting}

\subsection{2. Features to Monitor}

Monitor:

\begin{itemize}
    \item Load
    \item Lag 1
    \item Lag 24
    \item Lag 168
    \item Rolling mean 24
    \item Rolling mean 168
    \item Rolling standard deviation 24
    \item Hour
    \item Day of week
\end{itemize}

Calendar features are generally less informative for drift, but consumption-related features are useful signals.

\subsection{3. Evidently Integration}

Create:

\begin{lstlisting}
src/monitoring/drift.py
\end{lstlisting}

Save reports in:

\begin{lstlisting}
reports/drift/
\end{lstlisting}

For example:

\begin{lstlisting}
reports/drift/
+-- 2014_01.html
+-- 2014_02.html
+-- 2014_03.html
\end{lstlisting}

Store a machine-readable summary:

\begin{lstlisting}
{
    "dataset_drift": true,
    "drifted_features": 4,
    "total_features": 12,
    "drift_share": 0.333
}
\end{lstlisting}

\subsection{4. Retraining Condition}

Do not retrain simply because drift exists.

Use both data drift and performance degradation.

Recommended rule:

\begin{lstlisting}[language=Python]
should_retrain = (
    dataset_drift_detected
    and
    rolling_mape >
    production_baseline_mape * 1.25
)
\end{lstlisting}

This represents:

\[
\text{Retrain}
=
\text{Data Drift}
\land
\text{Performance Degradation}
\]

\subsection{5. Retraining Strategy}

Use an expanding training window.

Example:

Initial model:

\begin{lstlisting}
2012-01 through 2013-06
\end{lstlisting}

If retraining occurs in May 2014:

\begin{lstlisting}
Candidate training:
2012-01 through 2014-03

Candidate validation:
2014-04
\end{lstlisting}

Production resumes from May onward.

\subsection{6. Candidate Evaluation}

Compare:

\begin{itemize}
    \item Current production model
    \item Newly trained candidate
\end{itemize}

Use the same validation window.

Compare:

\begin{itemize}
    \item MAE
    \item RMSE
    \item MAPE
\end{itemize}

\subsection{7. Promotion Rules}

Do not deploy every retrained model.

Require measurable improvement.

Example:

\begin{lstlisting}[language=Python]
candidate_is_better = (
    candidate_mape <
    production_mape * 0.98
    and
    candidate_mae <
    production_mae
)
\end{lstlisting}

This requires approximately 2\% relative MAPE improvement.

Also impose an absolute quality requirement if desired:

\begin{lstlisting}[language=Python]
candidate_mape <= 0.08
\end{lstlisting}

\subsection{8. MLflow Retraining Records}

For each retraining event, log:

\begin{itemize}
    \item Trigger reason
    \item Drift result
    \item Production MAPE
    \item Training period
    \item Validation period
    \item Candidate metrics
    \item Production metrics
    \item Promotion decision
\end{itemize}

Example event:

\begin{lstlisting}
Run: retraining-2014-05

Trigger:
dataset_drift = true
rolling_mape = 7.2%

Production Model:
Version = 1
MAPE = 7.2%

Candidate Model:
Version = 2
MAPE = 5.8%

Decision:
PROMOTE
\end{lstlisting}

\subsection{9. Model Promotion}

The lifecycle becomes:

\begin{verbatim}
Production v1
     |
     v
Retrain
     |
     v
Candidate v2
     |
     v
Evaluation
     |
   Better?
   /    \
  No    Yes
  |      |
  v      v
Reject Promote
         |
         v
    Production v2
\end{verbatim}

\subsection{10. Rollback}

Implement a simple model-version selection mechanism.

Example conceptual functions:

\begin{lstlisting}[language=Python]
promote_model(version)
rollback_model(previous_version)
\end{lstlisting}

The goal is not to build a sophisticated deployment platform. The objective is to demonstrate awareness of model lifecycle and rollback.

\subsection{11. End-to-End Simulation}

The complete simulation should perform:

\begin{verbatim}
New Production Window
        |
        v
Generate Forecasts
        |
        v
Actual Data Arrives
        |
        v
Calculate Performance
        |
        v
Run Drift Detection
        |
        v
Performance Bad + Drift?
        |
      +---+---+
      |       |
     No      Yes
      |       |
      v       v
 Continue   Retrain
              |
              v
        Candidate Model
              |
              v
        Evaluate Candidate
              |
           Better?
          /      \
        No        Yes
        |          |
        v          v
      Reject     Promote
\end{verbatim}

\subsection{Day 6 Deliverables}

\begin{itemize}
    \item[\done] Evidently integration
    \item[\done] Reference/current monitoring windows
    \item[\done] Drift reports
    \item[\done] Retraining rules
    \item[\done] Candidate-model training
    \item[\done] Candidate evaluation
    \item[\done] Quality gates
    \item[\done] Model registration
    \item[\done] Model promotion
    \item[\done] Rollback support
\end{itemize}

\textbf{Stop condition:} A production simulation can automatically cause a new model version to become the active production model.

\newpage

% ==================================================
% DAY 7
% ==================================================

\section{Day 7 --- CI/CD, Testing, Documentation, and Portfolio Polish}

\subsection{Objective}

Turn the working technical system into a professional portfolio repository that can be understood quickly by recruiters and interviewers.

\subsection{1. GitHub Actions}

Create:

\begin{lstlisting}
.github/workflows/ci.yml
\end{lstlisting}

Trigger CI on:

\begin{itemize}
    \item Push
    \item Pull request
\end{itemize}

The CI pipeline should execute:

\begin{verbatim}
Checkout
   |
   v
Install Python
   |
   v
Install Dependencies
   |
   v
Ruff
   |
   v
pytest
   |
   v
Pipeline Sanity Check
   |
   v
Docker Build
\end{verbatim}

Avoid retraining the entire production model during every CI run.

Instead, use a small fixture dataset for automated tests.

\subsection{2. Test Suite}

Aim for approximately 15--25 meaningful tests.

\subsubsection{Data Tests}

\begin{itemize}
    \item \texttt{test\_no\_duplicates}
    \item \texttt{test\_hourly\_frequency}
    \item \texttt{test\_no\_negative\_load}
    \item \texttt{test\_no\_missing\_target}
\end{itemize}

\subsubsection{Feature Tests}

\begin{itemize}
    \item \texttt{test\_lag\_24}
    \item \texttt{test\_lag\_168}
    \item \texttt{test\_rolling\_features}
    \item \texttt{test\_no\_future\_leakage}
\end{itemize}

\subsubsection{Model Tests}

\begin{itemize}
    \item \texttt{test\_model\_training}
    \item \texttt{test\_prediction\_shape}
    \item \texttt{test\_prediction\_non\_negative}
    \item \texttt{test\_metrics\_generated}
\end{itemize}

\subsubsection{Monitoring Tests}

\begin{itemize}
    \item \texttt{test\_performance\_monitor}
    \item \texttt{test\_drift\_report\_created}
    \item \texttt{test\_retraining\_condition}
    \item \texttt{test\_no\_retraining\_when\_healthy}
\end{itemize}

\subsubsection{Registry Tests}

\begin{itemize}
    \item \texttt{test\_candidate\_comparison}
    \item \texttt{test\_promotion\_rule}
    \item \texttt{test\_failed\_candidate\_not\_promoted}
\end{itemize}

\subsubsection{API Tests}

\begin{itemize}
    \item \texttt{test\_health}
    \item \texttt{test\_model\_info}
    \item \texttt{test\_forecast}
    \item \texttt{test\_invalid\_input}
\end{itemize}

\subsection{3. Single-Command Demo}

Create:

\begin{lstlisting}[language=bash]
make demo
\end{lstlisting}

or:

\begin{lstlisting}[language=bash]
python scripts/run_demo.py
\end{lstlisting}

The demo should conceptually perform:

\begin{verbatim}
Preprocess
   |
   v
Feature Engineering
   |
   v
Initial Training
   |
   v
Register Model
   |
   v
Production Simulation
   |
   v
Monitor
   |
   v
Detect Drift
   |
   v
Retrain
   |
   v
Promote Candidate
   |
   v
Generate Reports
\end{verbatim}

\subsection{4. README Structure}

The README should contain:

\begin{enumerate}
    \item Overview
    \item Business / forecasting problem
    \item Dataset
    \item Architecture
    \item Forecasting approach
    \item Feature engineering
    \item MLOps pipeline
    \item Experiment tracking
    \item Drift monitoring
    \item Automated retraining
    \item Model promotion
    \item Results
    \item Local setup
    \item Docker instructions
    \item Testing
    \item Repository structure
    \item Engineering decisions
    \item Future improvements
\end{enumerate}

\subsection{5. README Architecture Diagram}

Include a diagram similar to:

\begin{verbatim}
                 UCI Electricity Data
                         |
                         v
                DVC Data Pipeline
                         |
                         v
                Feature Engineering
                         |
                         v
                    XGBoost
                         |
                         v
                MLflow Registry
                         |
                         v
                Forecast Service
                         |
               +---------+---------+
               |                   |
               v                   v
          Predictions            Actuals
               |                   |
               +---------+---------+
                         |
                         v
                    Monitoring
               +---------+---------+
               |                   |
               v                   v
            Drift            Forecast Error
               |                   |
               +---------+---------+
                         |
                         v
                 Retraining Gate
                         |
                         v
                  Candidate Model
                         |
                         v
                   Quality Gate
                    /        \
                   v          v
                Reject      Promote
\end{verbatim}

\subsection{6. Portfolio Screenshots}

Include screenshots of:

\begin{enumerate}
    \item MLflow experiment runs
    \item MLflow model registry
    \item Actual versus predicted forecast
    \item Evidently drift report
    \item FastAPI Swagger interface
    \item Passing GitHub Actions workflow
\end{enumerate}

\subsection{7. Report Real Results}

Avoid statements such as:

\begin{quote}
``The model performed well.''
\end{quote}

Instead report concrete results:

\begin{lstlisting}
Weekly Baseline MAPE: X%

Initial XGBoost MAPE: Y%

Production MAPE Before Retraining: Z%

Retrained Model MAPE: W%

Relative Improvement: N%
\end{lstlisting}

\subsection{8. Explain Engineering Decisions}

Document several key decisions.

\subsubsection{Why XGBoost?}

The objective is to demonstrate production ML lifecycle management rather than maximize forecasting benchmark performance.

A tree-based model:

\begin{itemize}
    \item Trains quickly
    \item Supports engineered lag features
    \item Is inexpensive to retrain
    \item Is simple to explain
    \item Keeps focus on MLOps
\end{itemize}

\subsubsection{Why Simulate Production?}

The source data is historical.

The 2014 period is therefore held out and replayed chronologically to simulate observations arriving in production.

\subsubsection{Why Require Both Drift and Performance Degradation?}

Data drift does not automatically imply that predictive performance has degraded.

Therefore:

\[
\text{Retrain}
=
\text{Drift Detected}
\land
\text{Performance Degraded}
\]

This avoids unnecessary retraining.

\subsection{Day 7 Deliverables}

\begin{itemize}
    \item[\done] GitHub Actions CI
    \item[\done] Complete test suite
    \item[\done] Docker build validation
    \item[\done] One-command demo
    \item[\done] Architecture diagram
    \item[\done] MLflow screenshots
    \item[\done] Drift-monitoring screenshots
    \item[\done] Forecast screenshots
    \item[\done] Strong README
    \item[\done] Documented engineering decisions
\end{itemize}

\textbf{Stop condition:} A recruiter should be able to understand the project within approximately two minutes of opening the repository.

\newpage

% ==================================================
% FINAL SYSTEM
% ==================================================

\section{Expected Final System Behavior}

The completed system should support the following workflow.

\subsection{Initial Training}

Run:

\begin{lstlisting}[language=bash]
dvc repro
\end{lstlisting}

The initial model is trained and registered.

Example:

\begin{lstlisting}
GridCastForecaster v1
Validation MAPE: 4.9%
Status: Production
\end{lstlisting}

\subsection{Production Simulation}

Run:

\begin{lstlisting}[language=bash]
python scripts/simulate_production.py
\end{lstlisting}

Expected output might resemble:

\begin{lstlisting}
[Week 01]
MAPE: 4.8%
Drift: False
Action: NONE

[Week 02]
MAPE: 5.0%
Drift: False
Action: NONE

[Week 03]
MAPE: 5.1%
Drift: False
Action: NONE

...

[Week 17]
MAPE: 6.7%
Drift: True
Action: MONITOR

[Week 18]
MAPE: 7.1%
Drift: True
Action: RETRAIN

Training candidate model...

Production model:
Version: 1
MAPE: 7.1%

Candidate model:
Version: 2
MAPE: 5.6%

Quality gate: PASSED

Promoting GridCastForecaster v2...

Current production model:
Version: 2
\end{lstlisting}

This sequence should become the main project demonstration.

\section{Recommended Daily Effort}

\begin{table}[h]
\centering
\renewcommand{\arraystretch}{1.2}
\begin{tabular}{clc}
\toprule
Day & Primary Focus & Approximate Time \\
\midrule
1 & Dataset, EDA, preprocessing & 6 hours \\
2 & Features and forecasting model & 7 hours \\
3 & MLflow and DVC & 7 hours \\
4 & FastAPI and Docker & 6 hours \\
5 & Production simulation and monitoring & 7 hours \\
6 & Drift, retraining, promotion & 8 hours \\
7 & CI/CD, tests, README, demo & 7 hours \\
\midrule
 & \textbf{Total} & \textbf{48 hours} \\
\bottomrule
\end{tabular}
\end{table}

\section{Project Priority}

If time becomes limited, prioritize features according to the following groups.

\subsection{Must Have}

\begin{itemize}
    \item Reproducible preprocessing
    \item Chronological train/validation split
    \item Baseline forecasting
    \item XGBoost model
    \item MLflow experiment tracking
    \item MLflow model versioning
    \item DVC pipeline
    \item Production-data replay
    \item Performance monitoring
    \item Drift detection
    \item Automated retraining
    \item Candidate evaluation
    \item Model promotion
    \item pytest
    \item Docker
    \item GitHub Actions
    \item Strong README
\end{itemize}

\subsection{Nice to Have}

\begin{itemize}
    \item FastAPI
    \item Automated HTML monitoring reports
    \item Makefile
    \item Rollback command
    \item Structured logging
    \item Feature-importance reports
\end{itemize}

\subsection{Skip During Week One}

Avoid:

\begin{itemize}
    \item Kubernetes
    \item Airflow
    \item Kafka
    \item Spark
    \item Terraform
    \item AWS infrastructure
    \item Prometheus
    \item Grafana
    \item Feature stores
    \item Microservices
    \item Deep neural networks
    \item Transformers
    \item Complex hyperparameter optimization
    \item Live streaming infrastructure
\end{itemize}

These technologies can be added later but are unnecessary for the initial portfolio version.

\section{Core MLOps Design Principle}

The project should clearly distinguish between three concepts.

\subsection{Data Drift}

Question:

\begin{quote}
Has the statistical distribution of production inputs changed?
\end{quote}

Conceptually:

\[
P_{\text{train}}(X)
\neq
P_{\text{production}}(X)
\]

\subsection{Performance Degradation}

Question:

\begin{quote}
Has the forecasting model become less accurate?
\end{quote}

For example:

\[
MAPE_{\text{production}}
>
MAPE_{\text{baseline}} \times 1.25
\]

\subsection{Retraining}

Question:

\begin{quote}
Can a newly trained model provide better performance than the currently deployed model?
\end{quote}

The final decision logic should therefore be:

\begin{verbatim}
Data Drift Detected
        +
Performance Degraded
        |
        v
Retrain Candidate
        |
        v
Evaluate Candidate
        |
        v
Compare with Production
        |
        v
Candidate Better?
     /       \
    No       Yes
    |         |
    v         v
 Reject    Promote
\end{verbatim}

A weak implementation would perform:

\begin{verbatim}
Drift Detected
     |
     v
Retrain
     |
     v
Automatically Deploy
\end{verbatim}

The stronger implementation used by \projectname{} is:

\begin{verbatim}
Drift Detected
      +
Performance Degraded
      |
      v
Retrain Candidate
      |
      v
Evaluate Candidate
      |
      v
Candidate Beats Production?
      |
   +--+--+
   |     |
  No    Yes
   |     |
   v     v
Reject Promote
\end{verbatim}

This distinction should be one of the main discussion points during interviews.

\section{Suggested Interview Explanation}

A concise explanation of the project can be:

\begin{quote}
I built an end-to-end electricity demand forecasting system that predicts aggregate power consumption using XGBoost and lag-based time-series features. The training pipeline is reproducible with DVC, experiments and model versions are tracked with MLflow, and the model is served through a containerized application.

I held out the final year of the dataset and replayed it chronologically to simulate production. During that simulation, the system monitors both feature drift and forecast error. If significant drift occurs together with performance degradation, the system automatically trains a candidate model. The candidate is evaluated against the existing production model and is promoted only if it passes predefined quality gates.
\end{quote}

\section{Definition of Done}

The project is complete when the following workflow can be demonstrated from the repository:

\begin{enumerate}
    \item Clone the repository
    \item Obtain the versioned dataset
    \item Run the reproducible DVC pipeline
    \item Train the initial model
    \item View experiments in MLflow
    \item Register the production model
    \item Start the forecasting service
    \item Replay production data
    \item Observe forecast performance
    \item Detect feature drift
    \item Trigger retraining
    \item Train a candidate model
    \item Evaluate candidate versus production
    \item Promote the candidate if it is better
    \item Continue forecasting with the new model version
    \item Verify the complete repository using automated tests and CI
\end{enumerate}

\vfill

\begin{center}
    \textbf{\projectname{} --- End-to-End Energy Forecasting MLOps Portfolio Project}
\end{center}

\end{document}